\section{Evaluation}
\label{sec:evaluation}

\subsection{Experimental Setup}
\label{sec:eval:setup}

We implement PFR using the \texttt{redist} Rust binary (\texttt{redist-apportion}
crate) with METIS $k$-way partitioning at each tree level. The per-level balance
tolerance is set to $\varepsilon_{\text{level}} = \varepsilon_{\text{statutory}} /
(d+1)$ where $d$ is the actual PFR tree depth for that state, so that the
cumulative imbalance across levels stays within the statutory $\varepsilon = 0.5\%$.
The METIS minimum ufactor of 0.1\% applies as a floor. A single
content-derived seed ($\text{SHA256}(\text{census\_release\_id} \|
\texttt{DIA\_SEED\_V1})$) is used throughout. Partisan outcomes are measured
against 2020 presidential two-party precinct returns interpolated to census tracts
\citep{Kuriwaki2023}. All 50 states at their 2020 apportionment seat counts are
evaluated.

\subsection{Factorization Structure of the 2020 Apportionment}
\label{sec:eval:factorizations}

Table~\ref{tab:factorizations} shows the 2020 seat count for each state and its
canonical PFR factorization (largest prime first). Seven seat counts are prime and
use binary fallback at their first level: CT and OK ($k=5$, $\to 2+3$),
SC and AL ($k=7$, $\to 3+4$), VA ($k=11$, $\to 5+6$), MI ($k=13$, $\to 6+7$),
and PA and IL ($k=17$, $\to 8+9$). Twelve states with $k = 2^m$ have fully binary
trees (MT, RI, NH, ME, HI, ID, WV: $k=2$; KS, UT, NV, AR, IA, MS: $k=4$;
CO, WI, MO, MD, MN: $k=8$) and are identical to standard recursive bisection.

\begin{table}[h]
\centering\small
\begin{tabular}{lrll|lrll}
\toprule
State & $k$ & Factorization & Depth &
State & $k$ & Factorization & Depth \\
\midrule
VT,AK,WY,DE  & 1  & ---          & 0 & WA & 10 & $5 \times 2$        & 3 \\
ND,SD         & 1  & ---          & 0 & VA & 11 & 11 (prime)$\to$5+6  & 3 \\
MT,RI,NH,ME  & 2  & $2$          & 1 & NJ & 12 & $3\times 2^2$       & 3 \\
HI,ID,WV     & 2  & $2$          & 1 & MI & 13 & 13 (prime)$\to$6+7  & 4 \\
NE,NM        & 3  & $3$          & 1 & GA & 14 & $7\times 2$         & 2 \\
KS,UT,NV,AR  & 4  & $2^2$        & 2 & NC & 14 & $7\times 2$         & 2 \\
IA,MS        & 4  & $2^2$        & 2 & OH & 15 & $5\times 3$         & 2 \\
CT,OK        & 5  & 5 (prime)$\to$2+3 & 2 & PA,IL & 17 & 17 (prime)$\to$8+9 & 4 \\
OR,KY,LA     & 6  & $3\times 2$  & 2 & NY & 26 & $13\times 2$        & 5 \\
SC,AL        & 7  & 7 (prime)$\to$3+4 & 2 & FL & 28 & $7\times 2^2$  & 3 \\
CO,WI,MO     & 8  & $2^3$        & 3 & TX & 38 & $19\times 2$        & 2 \\
MD,MN        & 8  & $2^3$        & 3 & CA & 52 & $13\times 2^2$      & 3 \\
TN,AZ,MA,IN  & 9  & $3^2$        & 2 & & & & \\
\bottomrule
\end{tabular}
\caption{2020 seat counts, PFR factorizations (largest prime first), and tree depth.
  Depth = number of METIS calls per branch from root to leaf.
  Prime seat counts (5, 7, 11, 13, 17, 19) use binary floor/ceil fallback at
  that level, which is counted as 1 METIS call.}
\label{tab:factorizations}
\end{table}

\subsection{National Comparison: PFR vs.\ MEC}
\label{sec:eval:national}

Table~\ref{tab:national-pfr-mec} compares the PFR partisan outcome to the
minimum-edge-cut (MEC) outcome from \citet{dellaLibera2026mec} for all 48
states with available precinct election data (Alaska and Hawaii excluded).

\begin{table}[h]
\centering\small
\begin{tabular}{lrrrr|lrrrr}
\toprule
State & $k$ & MEC\,D & PFR\,D & $\Delta$ &
State & $k$ & MEC\,D & PFR\,D & $\Delta$ \\
\midrule
GA & 14 & 7 & 5 & $-2$ & MN & 8  & 3  & 4  & $+1$ \\
VA & 11 & 8 & 6 & $-2$ & TN & 9  & 1  & 2  & $+1$ \\
OR &  6 & 4 & 3 & $-1$ & TX & 38 & 15 & 16 & $+1$ \\
SC &  7 & 3 & 2 & $-1$ & NC & 14 & 5  & 7  & $+2$ \\
NY & 26 & 21& 20& $-1$ & \multicolumn{5}{l}{\emph{34 states: identical outcome}} \\
PA & 17 & 8 & 7 & $-1$ & & & & & \\
WI &  8 & 3 & 2 & $-1$ & & & & & \\
NJ & 12 & 10& 9 & $-1$ & & & & & \\
\midrule
\multicolumn{5}{l|}{MEC national (48 states): 228D/204R (52.8\%\,D)} &
\multicolumn{5}{l}{PFR national (48 states): 223D/209R (51.6\%\,D)} \\
\bottomrule
\end{tabular}
\caption{PFR vs.\ MEC partisan outcomes. $\Delta$ = PFR\,D $-$ MEC\,D.
  34 states have identical outcomes; 12 differ by 1--2 seats. PFR gives
  Democrats 5 fewer seats nationally ($-2.3\%$ seat share).
  MEC outcomes from \citet{dellaLibera2026mec} at convergence
  (1400 seeds for WI; 1000 seeds for GA, TN, MN, PA).}
\label{tab:national-pfr-mec}
\end{table}

PFR produces 223D/209R vs.\ MEC's 228D/204R --- a 5-seat Republican shift.
Both totals are within 2 seats of proportional given the ~51\% aggregate
Democratic two-party vote across the 48 states. Neither algorithm targets
proportionality, and neither is systematically biased: the 5-seat gap reflects
the interaction of factorization-induced boundary structures with specific
state geographies, not a consistent partisan tilt of the algorithm.

\subsection{North Carolina and Georgia: Same Factorization, Opposite Effects}
\label{sec:eval:nc-ga}

The paper's central empirical finding concerns two states with identical
factorizations: both Georgia ($k=14$) and North Carolina ($k=14$) factor as
$14 = 7 \times 2$. Under PFR, both states receive the same hierarchical
prescription: a 7-way primary split of the state, then bisect each of the 7
sub-regions.

\begin{table}[h]
\centering
\begin{tabular}{lrrrrrr}
\toprule
State & $v_s$ & Factorization & PFR\,D & PFR\,$\Delta$ & MEC\,D & MEC\,$\Delta$ \\
\midrule
North Carolina & 49.3\% & $7\times 2$ & 7D/7R & $+0.7$\,pp & 5D/9R & $-13.6$\,pp \\
Georgia        & 50.1\% & $7\times 2$ & 5D/9R & $-14.4$\,pp & 7D/7R & $-0.1$\,pp \\
\bottomrule
\end{tabular}
\caption{North Carolina and Georgia: identical $k=14=7\times 2$ factorization,
  opposite partisan effects. $v_s$ = statewide Dem two-party vote share.
  $\Delta$ = proportionality gap ($+$ = Dem over-represented).
  NC under PFR achieves near-perfect proportionality ($+0.7$\,pp) while MEC
  produces $-13.6$\,pp. GA under PFR produces $-14.4$\,pp while MEC achieves
  $-0.1$\,pp.}
\label{tab:nc-ga}
\end{table}

For North Carolina, the 7-way primary split groups the state's Democratic
precincts (Charlotte, the Research Triangle, Greensboro) into configurations
that, when each sub-region is subsequently bisected, produce seven Democratic
and seven Republican districts --- near-perfect proportionality for a 49.3\%
Dem state. The same algorithm applied to Georgia, where Democratic voters
concentrate in Atlanta and its suburbs, produces a 5D/9R outcome that is
\emph{less} proportional than the MEC result.

This is not an algorithmic artifact. It is the core thesis of this paper:
\emph{the partition structure is determined by the prime factorization of the
seat count, and its partisan effect is determined entirely by the interaction of
that structure with the geographic distribution of voters}. The algorithm has no
partisan inputs and makes no partisan choices. The NC/GA divergence demonstrates
that the same factorization tree can produce outcomes that favour either party
depending on geography --- precisely the constitutional property we seek.

\paragraph{Wisconsin: single-seed limitation.}
Wisconsin ($k = 8 = 2^3$) is a fully binary PFR tree, identical in structure
to recursive bisection. PFR produces 2D/6R ($-25.3$\,pp) at the fixed seed.
This matches the 200-seed MEC \emph{false floor} from \citet{dellaLibera2026mec}:
the single PFR seed finds an early local optimum that MEC only escapes at seed
318 after 1{,}400 total seeds. For states with large near-minimum solution
spaces, a multi-seed PFR protocol (run several seeds, take the minimum-EC
factorization-tree result) would be needed to replicate MEC convergence quality.

\subsection{Prime Seat Counts: Pennsylvania and Michigan}
\label{sec:eval:prime}

Pennsylvania ($k = 17$, prime) and Michigan ($k = 13$, prime) use binary
floor/ceil fallback at their first split level.

For Pennsylvania, PFR produces 7D/10R ($-9.4$\,pp) vs.\ the MEC convergence
result of 8D/9R ($-3.5$\,pp). The PFR split at $k=17$ yields $\{8, 9\}$
sub-regions; the 8-district side ($2^3$, fully binary) and the 9-district side
($3^2$, fully ternary) find different local optima than the iterated bisection
that MEC uses. Pennsylvania's mixed geography (Philadelphia, Pittsburgh, Scranton,
broad rural areas) makes it sensitive to first-level boundary choices.

Michigan ($k = 13$, prime, $\to 6+7$) produces 7D/6R ($+2.4$\,pp) under both
PFR and MEC --- the outcomes agree despite the different algorithmic paths. The
6-district half ($2 \times 3$) and 7-district half (prime, $\to 3+4$) each
converge to the same partisan outcome as MEC's recursive bisection.

\subsection{Texas: Large-Prime Top-Level Split}
\label{sec:eval:texas}

Texas ($k = 38 = 2 \times 19$) has a depth-2 factorization tree: a 19-way
primary METIS split of the state, then bisect each of the 19 sub-regions.
PFR produces 16D/22R ($-5.1$\,pp) vs.\ MEC's 15D/23R ($-7.7$\,pp) --- one
more Democratic seat. The 19-way primary partition groups Texas's diverse urban
centres (Houston, Dallas, Austin, San Antonio) differently than iterated
bisection, yielding a slightly more proportional outcome. Texas also demonstrates
the computational tractability of large-prime top-level splits: a single METIS
19-way call at ufactor=1 completes in under one second per seed.

\subsection{Structural Disruption Across Reapportionments}
\label{sec:eval:disruption}

A quantitative disruption study (fraction of tracts reassigned when $k$ changes
by $\pm 1$ across 2010$\to$2020 reapportionments) is deferred to future work.
The key qualitative prediction is:

\begin{itemize}
  \item States that \emph{retain their largest prime factor}: share the canonical
        top-level partition across seat counts, so only lower levels of the tree
        change. Example: California $52 \to 53$ is a full structural break ($13$
        vs.\ prime), but if CA moved $52 \to 78 = 13 \times 6$, the 13-way
        primary partition would be reused.
  \item States that \emph{change their largest prime}: complete tree redesign.
        Example: $51 = 3 \times 17 \to 52 = 4 \times 13$ shares no top-level
        structure.
\end{itemize}

The cache in \texttt{PfrCompositor} directly implements this prediction:
computing $k=34$ and $k=51$ for the same state in sequence, the second call
records non-zero \texttt{cache\_hits} from the shared 17-way primary partition.
